
\documentclass[12pt,a4paper]{report}

% Packages
\usepackage[a4paper,margin=1in]{geometry}
\usepackage{xcolor}
\usepackage{sectsty}
\usepackage{fancyhdr}

\chapterfont{\color{red}\bfseries\Huge}
\sectionfont{\color{magenta}\itshape\Large}
\subsectionfont{\color{orange}\bfseries}

\pagestyle{fancy}
\fancyhf{}

\fancyhead[L]{Sanyam Sonwal}
\fancyhead[C]{Technical Writing}
\fancyhead[R]{\today}

\renewcommand{\headrulewidth}{0.4pt}
\setlength{\headheight}{15pt}

\begin{document}


\begin{titlepage}
    \thispagestyle{empty}
    \centering

    \vspace*{2cm}

    {\Huge\bfseries Introduction to LaTeX\par}

    \vspace{2cm}

    {\Large Sanyam Sonwal\par}

    \vspace{0.5cm}

    ID: 2025UCP1709

    \vspace{1cm}

    Course: Technical Writing

    \vspace{0.5cm}

    Course Instructor: Dr. Neeta Nain

    \vspace{1cm}

    Date: \today

    \vfill

    MNIT Jaipur

\end{titlepage}


\begin{abstract}
This document provides an introduction to LaTeX,
a document preparation system widely used for
technical and scientific writing. It explains the
basic structure of LaTeX documents, formatting
commands, chapters, sections, and other important
features used to create professional documents.
\end{abstract}

\chapter{Introduction}

LaTeX is a document preparation system used to
create professional and well-structured documents.
It is especially useful for scientific writing,
academic reports, research papers, and presentations.

\section{What is LaTeX?}

LaTeX is based on the TeX typesetting system.
It allows users to focus on content while maintaining
consistent document formatting.

\subsection{Applications of LaTeX}

LaTeX is used for academic reports, research papers,
books, and technical documentation.

\section{Features of LaTeX}

Some important features of LaTeX are:

\begin{itemize}
    \item Professional document formatting.
    \item Automatic numbering of chapters and sections.
    \item Mathematical equation support.
    \item Tables, figures, and references.
\end{itemize}

\chapter{Basic Structure of a LaTeX Document}

A LaTeX document consists of two main parts:
the preamble and the document environment.

\section{Preamble}

The preamble contains the document class,
packages, and formatting instructions.

\section{Document Environment}

The document environment contains the main
content of the report.

\chapter{Conclusion}

LaTeX is a powerful tool for creating structured
and professional documents. Its features make it
particularly useful for academic and technical writing.

\end{document}